% GENERATED FILE. Edit canonical lessons, then run npm run generate:books.
% canonical-source-hash: fnv1a64:64e77e56dc9cae17
% canonical-lessons: ML-C560-vinodayatra, ML-C560-citraracana, ML-C560-silpam, ML-C560-silpi, ML-C560-pradarsanam

\chapter{Sculptures and Exhibitions}
\label{ch:ml-sculptures-and-exhibitions}
\hlchaptermodality{560}

\begin{chapteropening}
\textbf{By the end of this chapter:} \emph{I can say an excursion, painting, a sculpture, a sculptor and an exhibition.}

Complete the last lesson of chapter 560: I can say an excursion, painting, a sculpture, a sculptor and an exhibition.
\end{chapteropening}

\section[\texorpdfstring{\ml{വിനോദയാത്ര} (vinōdayātra)}{വിനോദയാത്ര (vinōdayātra)}]{\ml{വിനോദയാത്ര} (vinōdayātra) — an excursion; a pleasure trip}

\label{lesson:ML-C560-vinodayatra}



\begin{quote}
Before the new one: say the Malayalam for to cherish, then the Malayalam for to ask forgiveness.
\end{quote}

\subsection*{You'll want to know: \ml{വിനോദയാത്ര}}
\textbf{\ml{വിനോദയാത്ര}} — \emph{vinōdayātra} — \textquotedblleft{}an excursion; a pleasure trip\textquotedblright{}.

\begin{quote}
\ml{സ്കൂളിൽ} \ml{നിന്ന്} \ml{കുട്ടികൾ} \ml{വിനോദയാത്രയ്ക്ക്} \ml{പോയി}.
\end{quote}

\begin{quote}
\emph{skūḷil ninnŭ kuṭṭikaḷ vinōdayātraykkŭ pōyi.}
\end{quote}

\begin{quote}
The children went on a school excursion.
\end{quote}

\ml{വിനോദം} a pastime + \ml{യാത്ര} a journey.

\begin{cousinweb}
\textbf{\ml{വിനോദയാത്ര}} (\emph{vinōdayātra}) is a learned compound built from Sanskrit \emph{vinoda}, \textquotedblleft{}diversion, pastime\textquotedblright{} (Monier-Williams), and \emph{yātrā}, \textquotedblleft{}a journey\textquotedblright{} (CDIAL 10456).
\end{cousinweb}

\begin{grammarlens}[title={What you've built}]
One more word for describing what you lived and felt.
\end{grammarlens}

\subsection*{Your turn}
\begin{itemize}
\raggedright
  \item \emph{Say it:} \emph{vinōdayātra}
  \item \emph{Say it:} \emph{vinōdayātra}, once more
  \item \emph{Say it:} \emph{skūḷil ninnŭ kuṭṭikaḷ vinōdayātraykkŭ pōyi.} — The children went on a school excursion.
  \item \emph{Recall:} say the Malayalam for to cherish, then the Malayalam for to ask forgiveness, then say \emph{vinōdayātra} again
\end{itemize}

\begin{tcolorbox}[breakable,colback=teal!4,colframe=teal!35!black,title={Before you move on}]
What does \ml{വിനോദയാത്ര} mean? (\textquotedblleft{}An excursion; a pleasure trip\textquotedblright{}.) Say it once more.
\end{tcolorbox}
\section[\texorpdfstring{\ml{ചിത്രരചന} (citraracana)}{ചിത്രരചന (citraracana)}]{\ml{ചിത്രരചന} (citraracana) — painting}

\label{lesson:ML-C560-citraracana}



\begin{quote}
Before the new one: say the Malayalam for to ask forgiveness, then the Malayalam for an excursion.
\end{quote}

\subsection*{You'll want to know: \ml{ചിത്രരചന}}
\textbf{\ml{ചിത്രരചന}} — \emph{citraracana} — \textquotedblleft{}painting\textquotedblright{}.

\begin{quote}
\ml{അനിയത്തിക്ക്} \ml{ചിത്രരചനയിൽ} \ml{സമ്മാനം} \ml{കിട്ടി}.
\end{quote}

\begin{quote}
\emph{aniyattikkŭ citraracanayil sammānaṁ kiṭṭi.}
\end{quote}

\begin{quote}
My little sister won a prize for painting.
\end{quote}

\ml{ചിത്രം} a picture + \ml{രചന} a making, a composition.

\begin{cousinweb}
\textbf{\ml{ചിത്രരചന}} (\emph{citraracana}) is a learned compound built from Sanskrit \emph{citra}, \textquotedblleft{}bright, variegated\textquotedblright{}, the source of \ml{ചിത്രം} (CDIAL 4803), and \emph{racana}, \textquotedblleft{}making, composing\textquotedblright{} (Monier-Williams).
\end{cousinweb}

\begin{grammarlens}[title={What you've built}]
One more word for describing what you lived and felt.
\end{grammarlens}

\subsection*{Your turn}
\begin{itemize}
\raggedright
  \item \emph{Say it:} \emph{citraracana}
  \item \emph{Say it:} \emph{citraracana}, once more
  \item \emph{Say it:} \emph{aniyattikkŭ citraracanayil sammānaṁ kiṭṭi.} — My little sister won a prize for painting.
  \item \emph{Recall:} say the Malayalam for to ask forgiveness, then the Malayalam for an excursion, then say \emph{citraracana} again
\end{itemize}

\begin{tcolorbox}[breakable,colback=teal!4,colframe=teal!35!black,title={Before you move on}]
What does \ml{ചിത്രരചന} mean? (\textquotedblleft{}Painting\textquotedblright{}.) Say it once more.
\end{tcolorbox}
\section[\texorpdfstring{\ml{ശില്പം} (śilpaṁ)}{ശില്പം (śilpaṁ)}]{\ml{ശില്പം} (śilpaṁ) — a sculpture; a statue}

\label{lesson:ML-C560-silpam}



\begin{quote}
Before the new one: say the Malayalam for an excursion, then the Malayalam for painting.
\end{quote}

\subsection*{You'll want to know: \ml{ശില്പം}}
\textbf{\ml{ശില്പം}} — \emph{śilpaṁ} — \textquotedblleft{}a sculpture; a statue\textquotedblright{}.

\begin{quote}
\ml{ക്ഷേത്രത്തിലെ} \ml{കല്ലുകൊണ്ടുള്ള} \ml{ശില്പങ്ങൾ} \ml{വളരെ} \ml{പഴയതാണ്}.
\end{quote}

\begin{quote}
\emph{kṣētrattile kallukoṇṭuḷḷa śilpaṅṅaḷ vaḷare paḻayatāṇŭ.}
\end{quote}

\begin{quote}
The stone sculptures of the temple are very old.
\end{quote}

Plural: \ml{ശില്പങ്ങൾ}, \emph{śilpaṅṅaḷ}.

Plural: \ml{ശില്പങ്ങൾ}.

\begin{cousinweb}
\textbf{\ml{ശില്പം}} (\emph{śilpaṁ}) is borrowed from Sanskrit \emph{śilpa}, \textquotedblleft{}artistic work, a manual craft\textquotedblright{} (CDIAL 12470).
\end{cousinweb}

\begin{grammarlens}[title={What you've built}]
One more word for describing what you lived and felt.
\end{grammarlens}

\subsection*{Your turn}
\begin{itemize}
\raggedright
  \item \emph{Say it:} \emph{śilpaṁ}
  \item \emph{Say it:} \emph{śilpaṁ}, once more
  \item \emph{Say it:} \emph{kṣētrattile kallukoṇṭuḷḷa śilpaṅṅaḷ vaḷare paḻayatāṇŭ.} — The stone sculptures of the temple are very old.
  \item \emph{Recall:} say the Malayalam for an excursion, then the Malayalam for painting, then say \emph{śilpaṁ} again
\end{itemize}

\begin{tcolorbox}[breakable,colback=teal!4,colframe=teal!35!black,title={Before you move on}]
What does \ml{ശില്പം} mean? (\textquotedblleft{}A sculpture; a statue\textquotedblright{}.) Say it once more.
\end{tcolorbox}
\section[\texorpdfstring{\ml{ശില്പി} (śilpi)}{ശില്പി (śilpi)}]{\ml{ശില്പി} (śilpi) — a sculptor}

\label{lesson:ML-C560-silpi}



\begin{quote}
Before the new one: say the Malayalam for painting, then the Malayalam for a sculpture.
\end{quote}

\subsection*{You'll want to know: \ml{ശില്പി}}
\textbf{\ml{ശില്പി}} — \emph{śilpi} — \textquotedblleft{}a sculptor\textquotedblright{}.

\begin{quote}
\ml{ആ} \ml{ശില്പി} \ml{തടിയിൽ} \ml{ഒരു} \ml{ആനയെ} \ml{കൊത്തി}.
\end{quote}

\begin{quote}
\emph{ā śilpi taṭiyil oru ānaye kotti.}
\end{quote}

\begin{quote}
That sculptor carved an elephant in wood.
\end{quote}

One who makes \ml{ശില്പങ്ങൾ}; \ml{വാസ്തുശില്പി}, an architect, holds it.

\begin{cousinweb}
\textbf{\ml{ശില്പി}} (\emph{śilpi}) is borrowed from Sanskrit \emph{śilpin}, \textquotedblleft{}skilled in art, an artisan\textquotedblright{} (CDIAL 12471).
\end{cousinweb}

\begin{grammarlens}[title={What you've built}]
One more word for describing what you lived and felt.
\end{grammarlens}

\subsection*{Your turn}
\begin{itemize}
\raggedright
  \item \emph{Say it:} \emph{śilpi}
  \item \emph{Say it:} \emph{śilpi}, once more
  \item \emph{Say it:} \emph{ā śilpi taṭiyil oru ānaye kotti.} — That sculptor carved an elephant in wood.
  \item \emph{Recall:} say the Malayalam for painting, then the Malayalam for a sculpture, then say \emph{śilpi} again
\end{itemize}

\begin{tcolorbox}[breakable,colback=teal!4,colframe=teal!35!black,title={Before you move on}]
What does \ml{ശില്പി} mean? (\textquotedblleft{}A sculptor\textquotedblright{}.) Say it once more.
\end{tcolorbox}
\section[\texorpdfstring{\ml{പ്രദർശനം} (pradarśanaṁ)}{പ്രദർശനം (pradarśanaṁ)}]{\ml{പ്രദർശനം} (pradarśanaṁ) — an exhibition; a show}

\label{lesson:ML-C560-pradarsanam}



\begin{quote}
Before the new one: say the Malayalam for a sculpture, then the Malayalam for a sculptor.
\end{quote}

\subsection*{You'll want to know: \ml{പ്രദർശനം}}
\textbf{\ml{പ്രദർശനം}} — \emph{pradarśanaṁ} — \textquotedblleft{}an exhibition; a show\textquotedblright{}.

\begin{quote}
\ml{നഗരത്തിൽ} \ml{ഒരു} \ml{പുസ്തക} \ml{പ്രദർശനം} \ml{നടക്കുന്നു}.
\end{quote}

\begin{quote}
\emph{nagarattil oru pustaka pradarśanaṁ naṭakkunnu.}
\end{quote}

\begin{quote}
A book exhibition is on in the city.
\end{quote}

Also a screening of a film: \ml{സിനിമയുടെ} \ml{പ്രദർശനം}.

\begin{cousinweb}
\textbf{\ml{പ്രദർശനം}} (\emph{pradarśanaṁ}) is borrowed from Sanskrit \emph{pradarśana}, \textquotedblleft{}appearance, showing, pointing out\textquotedblright{} (Monier-Williams).
\end{cousinweb}

\begin{grammarlens}[title={What you've built}]
One more word for describing what you lived and felt.
\end{grammarlens}

\subsection*{Your turn}
\begin{itemize}
\raggedright
  \item \emph{Say it:} \emph{pradarśanaṁ}
  \item \emph{Say it:} \emph{pradarśanaṁ}, once more
  \item \emph{Say it:} \emph{nagarattil oru pustaka pradarśanaṁ naṭakkunnu.} — A book exhibition is on in the city.
  \item \emph{Recall:} say the Malayalam for a sculpture, then the Malayalam for a sculptor, then say \emph{pradarśanaṁ} again
\end{itemize}

\begin{tcolorbox}[breakable,colback=teal!4,colframe=teal!35!black,title={Before you move on}]
What does \ml{പ്രദർശനം} mean? (\textquotedblleft{}An exhibition; a show\textquotedblright{}.) Say it once more.
\end{tcolorbox}
